% ---
% id: protocol-compiler/def:homog
% kind: def
% title: "Licensed constraint list"
% summary: >-
%   A constraint list is licensed when every entry is declared in the form its
%   source mode admits, homogeneous scalar data in trusted mode and scalars in
%   untrusted mode, or carries the party-local tag. Licensing reads the
%   declared tag and compares no numbers.
% prominence: standard
% status: implemented
% uses: [def:declaration, def:feasible]
% mentions: [thm:gates, reductions/def:licensed, reductions/cor:entries, reductions/req:party-local-tag, foundations/req:consumer-decides-exactly, req:no-window]
% assumes: []
% refs: []
% provenance: newly-derived
% implements: certiqs_key.protocol_compiler.gates:licensed_for_the_emitted_mode
% lean:
% notes: "reductions/def:licensed and reductions/cor:entries are mentions and not uses: this definition states the licensed class in its own words, and a uses edge would carry the soundness ledger of the chain's producer into every compiler theorem."
% ---

Reading the tag of an entry off the type of its datum is a decision on the
data handed over, and infers nothing from a property of an object the entry
was built from (\xref{foundations/req:consumer-decides-exactly}). The
reductions module, which defines the \tref{reduction chain} this module consumes,
declares every entry's right-hand side as a real constant
(\xref{reductions/def:licensed}, \xref{reductions/cor:entries}). On a
one-factor emission of that module the form beside it is the tag value
homogeneous in trusted-source mode, so every entry is homogeneous, the
licensed form in trusted-source mode, and an arriving tag field agrees with
the type read off the datum; where the \tref{measured system} is a declared product
of party-scoped factors, an entry may instead arrive tagged party-local, with
a scope (\xref{reductions/req:party-local-tag}). That licensing is
decided exactly and at no window is part of Requirement~\ref{req:no-window}.

\begin{definition}\label{def:homog}
Each entry of $\CList$ carries its right-hand side in a declared form, and
the form is part of the datum rather than a property measured of it
(Definition~\ref{def:declaration}). There are two \emph{untagged} forms,
called so because their tag is the type of their datum and needs no field of
its own: a
\term[Said of a constraint entry, not declared party-local, whose right-hand side is a real scalar; in trusted-source mode it denotes that scalar times the marginal.]{homogeneous}
entry declares a real scalar $c_k$, a
\term[Said of a constraint entry declaring a Hermitian matrix on Alice's system as its right-hand side.]{general}
entry declares a Hermitian $\Theta_k$ on $A$; in untrusted-source mode only
the scalar form exists. For these two forms
the tag \emph{is} the type of the right-hand side an entry carries, and it is
read off that type: an entry whose right-hand-side slot holds a real scalar
is homogeneous, and one whose slot holds a Hermitian matrix on $A$ is
general. \emph{That
reading is a type decision}: it inspects which of two objects a slot holds,
compares no numbers, reads no matrix entry and measures nothing. The third
form, party-local, is declared by its tag and is not read off the type of the
datum. $\CList$ is
\term[Said of a constraint list every entry of which is in the form the declared source mode admits, or is party-local.]{licensed for the declared source mode}
when every entry is in the form that mode admits:
\begin{enumerate}[label=(\alph*),leftmargin=2.2em]
\item \emph{trusted-source mode:} every entry is declared
  homogeneous, carrying $c_k\in\mathbb{R}$; the right-hand side it denotes is
  then \emph{constructed} as $\Theta_k\coloneqq c_k\tau_A$ from the declared
  $c_k$ and the declared $\tau_A$, and the entry of \eqref{eq:feas-modes}
  reads $\tr_{\Qsys}[(\idon{A}\otimes\hat O_k)\,\omega]\bowtie_kc_k\tau_A$;
\item \emph{untrusted-source mode:} every right-hand side is a real number
  $\Theta_k=c_k$ and the entry reads
  $\tr[(\idon{A}\otimes\hat O_k)\,\omega]\bowtie_kc_k$.
\end{enumerate}
Licensing is therefore a \emph{type} check on the tag, and it compares no
numbers: what is inspected is which of the two forms an entry declares; no
operator identity is tested against any number, and nothing here is decided
at a window. \emph{Only in cases (a) and
(b) are the entries written $(\hat O_k,\bowtie_k,c_k)$}; a
general entry is written $(\hat O_k,\bowtie_k,\Theta_k)$ and is the form the
notation table records under $\Theta_k$. A general entry is admissible
input, every structural claim of this contract being proved for it, but it is
not licensed, so a list containing one carries no key-length claim
(Theorem~\ref{thm:gates}).

\emph{The \tref{party-local form}, a third licensed form.} In the general format of
Definition~\ref{def:declaration}\,(DX6) an entry may be declared
\term[Said of a scoped constraint entry, tagged so, whose condition holds against the marginal of the state on the systems other than its factor.]{party-local}:
it carries its scope $f$, a block list $\hat O_k$ against the
\tref{sector map} of the \tref{measured factor} $f$ and a real $c_k$, it is
written
$(\hat O_k,\bowtie_k,c_k;f)$, and its state-level reading, in either source
mode, is
\begin{equation}\label{eq:partylocal-entry}
  \tr_{\Qsys_f}\!\Big[\big(\idon{}\otimes
  (\hat O_k-c_k\,\idon{\Qsys_f})\big)\,\omega\Big]\;\bowtie_k\;0
  \qquad\text{as an operator on the systems of }\omega
  \text{ other than }\Qsys_f ,
\end{equation}
those systems being $A$ and the measured factors $\Qsys_{f'}$, $f'\ne f$;
equivalently $\tr_{\Qsys_f}[(\idon{}\otimes\hat O_k)\,\omega]\bowtie_k
c_k\tr_{\Qsys_f}\omega$, the entry being homogeneous with respect to the
marginal that is actually there. Block by block it is the last line of
\eqref{eq:feas-general}. The list is then \emph{licensed} when every entry
is either in the form the source mode admits, (a) or (b), read for its
\tref{lifted operator}, or is declared party-local. This is still a type
check and compares no numbers: what is inspected is the declared tag. A
party-local entry's right-hand-side slot holds a real scalar, as a
homogeneous entry's does, so the type of the datum does not separate the
two: an entry is party-local exactly when it carries that tag.

\emph{The soundness condition of a party-local entry.} A party-local entry
$(\hat O_k,\bowtie_k,c_k;f)$ meets the
\term[The condition a party-local entry is sound under: every attack reaches its scoped factor through a channel on that factor alone whose adjoint, applied to the entry's operator, stands in the entry's relation to its scalar times the identity.]{soundness condition of a party-local entry}
when every attack reaches $\Qsys_f$ through a channel $\Lambda_f$ acting on
that factor alone which carries the spectral certificate
$\Lambda_f^{\dagger}(\hat O_k)\bowtie_kc_k\idon{}$, that is, the adjoint of
the channel applied to the entry's operator stands in the entry's relation to
$c_k$ times the identity.

\emph{What the tag changes, case by case.} With one measured factor and
trusted-source mode, \eqref{eq:partylocal-entry} reads
$\tr_{\Qsys}[(\idon{A}\otimes\hat O_k)\,\omega]\bowtie_kc_k\,\omega_A$, which
under the marginal condition $\omega_A=\tau_A$ is the form (a) of the same
$(\hat O_k,\bowtie_k,c_k)$. With one measured factor and $d_A=1$, in either
mode, it is the comparison of real numbers
$\tr[\hat O_k\omega]\bowtie_kc_k\tr\omega=c_k$, the form (b). \emph{With one
measured factor, untrusted-source mode and $d_A>1$ it is neither}: it is an
operator inequality on $A$ against the marginal $\omega_A$ that is actually
there, and the scalar entry (b) of the same $(\hat O_k,\bowtie_k,c_k)$ is its
trace, implied by it and not equivalent to it. With more than one measured
factor it implies, in either mode, the untagged reading of the same entry,
which is obtained from it by the partial trace over the other measured
factors (and over $A$ in untrusted-source mode), and it is in general not
implied by it. So the tag never weakens an entry, and of the four one-factor
cases, by source mode and by whether $d_A=1$, it changes something only in
untrusted-source mode at $d_A>1$.

\emph{Which relation symbols a tagged entry carries, in either mode.} An entry
declared party-local carries $\bowtie_k\in\{\preceq,=,\succeq\}$ \emph{in
either source mode}, because \eqref{eq:partylocal-entry} is a relation
between \emph{operators} in either mode: its left-hand side is an operator on
the systems of $\omega$ other than $\Qsys_{f_k}$, which are $A$ and the other
measured factors. The scalar case is the degenerate one: where
$d_{A_t}\prod_{f\ne f_k}d_{f,s_f}=1$ for every row --- every Alice block, and
every sector block of every measured factor other than $f_k$, of dimension
one, which with one measured factor means every Alice block of dimension one,
whether $d_A=1$ or $d_A=|\mathbb{T}|>1$ --- each row is a comparison of
$1\times1$ matrices, that is of real numbers, and the three operator symbols
are then the three scalar symbols $\le$, $=$, $\ge$. With one measured factor
and $d_A=1$ there is one row, and the reading is the form (b) of the same
$(\hat O_k,\bowtie_k,c_k)$, as the case-by-case paragraph states; otherwise
there is one such comparison per row, and the reading relates to the untagged
forms as that paragraph states for its case. \emph{Nothing about the
untagged forms changes}: an untagged entry carries the symbols its mode
admits, and the operator symbols beside a tag are a property of the third
form and of no other. No entry of a declaration in the \tref{basic format} is
tagged.
\end{definition}
